\documentclass[../RFC-HDFG-2026-003.tex]{subfiles}

\begin{document}

\section{Introduction}
\label{sec:intro}

\subsection{Purpose}

This document describes the design for in-file blob configuration
storage for HDF5 I/O filters: a binary channel for filter configuration
data too large for the \texttt{cd\_values} array or the parameter-string
API introduced by RFC-HDFG-2026-001 (string-based filter configuration).

\subsection{Scope}

This design covers:

\begin{itemize}
\item Three new, optional callbacks on \texttt{H5Z\_class3\_t}:
  \texttt{write\_blob}, \texttt{read\_blob}, and \texttt{close\_blob}.
\item A new public API function, \texttt{H5Pappend\_filter\_blob}.
\item A pipeline-message on-disk format extension to carry a blob
  alongside a filter's existing \texttt{cd\_values}.
\item Parallel I/O, \texttt{H5Pencode}/\texttt{H5Pdecode}, and
  \texttt{h5repack} semantics for blob-configured filters.
\end{itemize}

This design does not cover:

\begin{itemize}
\item Changes to \texttt{H5Z\_func\_t}, the runtime I/O callback used by
  \texttt{H5Z\_class1\_t}/\texttt{H5Z\_class2\_t} filters.
\item Changes to \texttt{H5Z\_class2\_t} or \texttt{H5Z\_class1\_t}.
\item Blob support for group fractal heap filters.
\item Multi-file blob sharing (a blob in file~A referenced by a dataset
  in file~B).
\end{itemize}

\subsection{Relationship to RFC-HDFG-2026-001}
\label{sec:relationship}

This design was, for a short period, folded directly into
RFC-HDFG-2026-001 as a single feature before being split back out into
this standalone RFC; the resolved design decisions in
\SectionRef{sec:decisions} carry forward unchanged from that period.

RFC-HDFG-2026-001 does not reserve any placeholder fields in
\texttt{H5Z\_class3\_t} for this RFC. An earlier draft reserved three
\texttt{void~*} fields (\texttt{write\_blob}/\texttt{read\_blob}/\texttt{close\_blob})
for later activation by this RFC; that approach was dropped because
pre-reserving an opaque placeholder does not avoid the struct-version
bump this feature needs once real callback signatures are decided --- it
only defers the decision. Consequently, this RFC's relationship to
\texttt{H5Z\_class3\_t} depends on timing:

\begin{itemize}
\item \textbf{If RFC-HDFG-2026-001 has not yet shipped} when this RFC is
  adopted, the three blob callbacks (\SectionRef{sec:blob-callbacks}) are
  added directly to \texttt{H5Z\_class3\_t}, exactly as if they had been
  part of RFC-HDFG-2026-001 from the start. No new struct version is
  required.
\item \textbf{If \texttt{H5Z\_class3\_t} has already shipped}, adding
  these fields requires a new \texttt{H5Z\_class4\_t} (or equivalent
  extension mechanism, e.g. the \texttt{H5Z\_CLASS\_T\_VERS\_MAX} pattern
  documented in RFC-HDFG-2026-001). This RFC does not resolve that
  question in advance; whoever finalizes this RFC should check the
  shipped status of \texttt{H5Z\_class3\_t} at that time and adjust
  accordingly.
\end{itemize}

\subsection{Motivation}
\label{sec:motivation}

RFC-HDFG-2026-001 raised the practical filter-configuration-size limit
from 256 \texttt{cd\_values} elements (~1~KB) to 4096 characters via a
human-readable parameter string. This is sufficient for most filters
but not for all. Two use cases from LibPressio (raised by @robertu94
during the community review of RFC-HDFG-2026-001,
HDFGroup/hdf5\_plugins\#255; see \texttt{LIBPRESSIO-FEEDBACK.md} in the
RFC-HDFG-2026-001 repository) remain unaddressed:

\verysubsection{Use Case A --- SZ4 JIT compiler (primary; fundamental size limit)}
SZ4 features a just-in-time compiler for fast user-defined modules. When
used, it needs to store pre-processed source code with the dataset ---
potentially multi-megabyte. There is no external dataset to reference;
the binary blob must travel with the dataset and be recovered on open.
This is a pure size limitation that \texttt{cd\_values} and the
parameter string cannot address.

\verysubsection{Use Case B --- ROIBIN-SZ spatial mask (secondary; reference to another dataset)}
ROIBIN-SZ stores a binary mask indicating which values use lossless
compression. The preferred design is for this mask to live as a
user-visible HDF5 dataset, with only a reference (a small, fixed-size
byte sequence) stored in the filter configuration. The filter
dereferences the reference at open time to recover the mask.

Both cases are solved by the same mechanism: a binary blob, opaque to
the library, stored alongside the dataset's pipeline message and handed
back to the filter verbatim at open time. This RFC specifies that
mechanism as a single, complete feature.

\subsection{Goals}

\begin{enumerate}
\item Filters can store configuration data of essentially arbitrary size,
  not bounded by \texttt{cd\_values} or the parameter-string limit.
\item Filters that do not need custom on-disk layout use a default
  storage implementation with no additional code.
\item The feature is correct under parallel I/O, DCPL
  encode/decode, and \texttt{h5repack} cross-file copy.
\item The runtime chunk I/O hot path is unaffected.
\end{enumerate}

\end{document}
